\doublespacing % Do not change - required
\label{Appendix}
\chapter{Experimental Configuration}

Table~\ref{tab:app-reproduction} summarises the retained hardware,
calibration and evaluation settings.

\begin{table}[htbp]
\centering
\caption{Experimental configuration used for the retained experiments.}
\label{tab:app-reproduction}

\begin{spacing}{1}
\small
\renewcommand{\arraystretch}{1.18}

\begin{tabularx}{\textwidth}{@{}
>{\raggedright\arraybackslash}p{0.40\textwidth}
>{\raggedright\arraybackslash}X
@{}}
\toprule
\textbf{Setting} & \textbf{Retained configuration} \\
\midrule

Source model &
\texttt{meta-llama/Llama-3.2-3B-Instruct} \\

Source precision &
BF16 \\

GPU &
NVIDIA GeForce RTX~5090 \\

CUDA version &
12.8 \\

PyTorch version &
2.11.0 \\

Transformers version &
5.10.1 \\

Decoder layers &
28 \\

Target Linear modules &
196 \\

Target weight elements &
2,818,572,288 \\

Calibration dataset &
UltraChat-200k \texttt{train\_sft} \\

NestQuant calibration context &
2,048 tokens \\

Evaluation dataset &
UltraChat-200k \texttt{test\_sft} \\

Evaluation target &
Assistant-response tokens only \\

Assistant responses &
73,150 \\

Scored assistant tokens &
4,500,809 \\

Maximum sequence length &
256 tokens \\

Maximum scored response length &
64 tokens \\

Batch size &
1 \\

Warm-up &
3 examples excluded from measurement \\

Primary quality metric &
Perplexity \\

Cross-method quality reference &
BF16 \\

NestQuant scope-ablation reference &
W4.25A16KV16 \\

\bottomrule
\end{tabularx}
\end{spacing}
\end{table}

Across the 28 decoder layers, this corresponds to 196 target Linear modules and
2,818,572,288 target weight elements. Embeddings, normalisation
parameters, the language-model head and other non-target tensors remain
outside this common Linear-layer quantization scope unless explicitly
stated otherwise.

All principal experiments use
\texttt{meta-llama/Llama-3.2-3B-Instruct} as the common source model.
The source checkpoint is represented using BF16 parameters and provides
the full-precision reference for the cross-method comparisons. The main
experiments were executed on a single NVIDIA GeForce RTX~5090 using
CUDA~12.8, PyTorch~2.11.0 and Transformers~5.10.1.

\chapter{Implementation Details}

The project-specific W4A8 pipeline requires a calibration stage to
select the scaling transformation used before final weight
quantization. Algorithm~\ref{alg:w4a8-calibration} summarises this
procedure. Candidate scales are evaluated using block-output
reconstruction error, after which the retained weights are quantized
using grouped symmetric INT4 and stored in packed form.

\begin{algorithm}[htbp]
\caption{Project-specific W4A8 calibration and weight construction}
\label{alg:w4a8-calibration}

\KwIn{
BF16 model $M$;
calibration set $\mathcal{D}_{\mathrm{cal}}$;
target Linear layers $\mathcal{L}$;
group size $G=128$;
candidate scale set $\mathcal{S}=\{s^{(1)},\ldots,s^{(20)}\}$
}

\KwOut{
Packed INT4 weights and retained scaling parameters
}

Collect representative input activations for all
$\ell \in \mathcal{L}$ using $\mathcal{D}_{\mathrm{cal}}$\;

\ForEach{target Linear layer $\ell \in \mathcal{L}$}{

    Compute the BF16 reference block output
    $Y_{\mathrm{BF16}}$\;

    $E_{\min} \leftarrow +\infty$\;
    $s^{*} \leftarrow \varnothing$\;

    \ForEach{candidate scale $s \in \mathcal{S}$}{

        Apply the candidate channel-wise scaling transformation\;

        Quantize the transformed weights using grouped symmetric
        INT4 with group size $G$\;

        Simulate the intended W4A8 numerical path to obtain
        $Y_{\mathrm{W4A8}}(s)$\;

        $E(s) \leftarrow
        \frac{1}{N}
        \sum_{n=1}^{N}
        \left\|
        Y_{\mathrm{BF16}}^{(n)}
        -
        Y_{\mathrm{W4A8}}^{(n)}(s)
        \right\|_2^2$\;

        \If{$E(s)<E_{\min}$}{
            $E_{\min} \leftarrow E(s)$\;
            $s^{*} \leftarrow s$\;
        }
    }

    Reconstruct the layer using the retained scale $s^{*}$\;

    Quantize the final weights using symmetric grouped INT4\;

    Pack the INT4 weight codes for runtime execution\;
}

\Return packed W4 weights and associated scaling parameters\;

\end{algorithm}

After calibration, the W4A8 configuration combines the retained packed
4-bit weights with dynamically quantized 8-bit activations during
inference. Algorithm~\ref{alg:w4a8-runtime} describes the runtime
activation-quantization procedure and the subsequent low-precision
Linear operation.

\begin{algorithm}[htbp]
\caption{Dynamic W4A8 runtime execution for one Linear layer}
\label{alg:w4a8-runtime}

\KwIn{
Floating-point activation matrix $X$;
packed INT4 weights $W_{4}$;
weight scales $S_W$
}

\KwOut{
Higher-precision output activation $Y$
}

\ForEach{input token row $x_i$ in $X$}{

    $s_{x_i}
    \leftarrow
    \dfrac{\max_j |x_{i,j}|}{127}$\;

    $q_{x_i}
    \leftarrow
    \operatorname{clip}
    \left(
    \operatorname{round}
    \left(
    \dfrac{x_i}{s_{x_i}}
    \right),
    -127,127
    \right)$\;
}

$X_{8} \leftarrow
\operatorname{stack}(q_{x_1},\ldots,q_{x_m})$\;

Execute the low-bit Linear path using
$X_{8}$, packed $W_{4}$, $S_W$ and the activation scales\;

Accumulate the matrix product in the supported higher-precision
accumulator\;

Apply the corresponding dequantization/scaling factors\;

$Y \leftarrow$ reconstructed output activation\;

\Return $Y$\;

\end{algorithm}

NestQuant differs from the scalar and group-wise integer quantizers by
representing eight-dimensional weight vectors using a structured
\(E_8\) lattice. Algorithm~\ref{alg:nestquant-packing} summarises the
selection of the candidate scaling coefficient and the construction of
the retained packed representation.

\begin{algorithm}[htbp]
\caption{NestQuant vector quantization and W4.25 packing}
\label{alg:nestquant-packing}

\KwIn{
Weight vector block $v\in\mathbb{R}^{8}$;
candidate scales
$\{\beta_1,\ldots,\beta_k\}$ with $k=4$;
$E_8$ lattice quantizer
}

\KwOut{
Packed lattice code and 2-bit scale index
}

Apply the retained normalisation procedure to $v$\;

Apply the orthogonal Hadamard rotation\;

$D_{\min} \leftarrow +\infty$\;
$p^{*} \leftarrow 0$\;
$c^{*} \leftarrow \varnothing$\;

\For{$p\leftarrow1$ \KwTo $k$}{

    $\tilde{v}
    \leftarrow
    \dfrac{v}{\beta_p}$\;

    $c_p
    \leftarrow
    \operatorname{Encode}_{E_8}(\tilde{v})$\;

    $\hat{v}_p
    \leftarrow
    \beta_p
    \operatorname{Decode}_{E_8}(c_p)$\;

    $D_p
    \leftarrow
    \left\|
    v-\hat{v}_p
    \right\|_2^2$\;

    \If{$D_p<D_{\min}$}{
        $D_{\min}\leftarrow D_p$\;
        $p^{*}\leftarrow p$\;
        $c^{*}\leftarrow c_p$\;
    }
}

Store each coordinate code of $c^{*}$ using a fixed 4-bit field\;

Store $p^{*}$ using a 2-bit scale index\;

$R_{\mathrm{packed}}
\leftarrow
\dfrac{8\times4+2}{8}
=
4.25$ bits per weight\;

\Return $(c^{*},p^{*})$\;

\end{algorithm}




\chapter{Supplementary Results}
\label{app:supplementary-results}

This appendix consolidates the numerical results underlying the
principal comparisons in Chapter~4. The tables provide a compact
numerical counterpart to the figures presented in the main text.

Table~\ref{tab:app-quality-results} reports the complete cross-method
perplexity comparison relative to the common BF16 baseline.

\begin{table}[htbp]
\centering
\caption{Complete cross-method model-quality results.}
\label{tab:app-quality-results}

\begin{spacing}{1}
\small
\renewcommand{\arraystretch}{1.18}

\begin{tabularx}{0.92\textwidth}{@{}
>{\raggedright\arraybackslash}X
>{\centering\arraybackslash}p{0.18\textwidth}
>{\centering\arraybackslash}p{0.20\textwidth}
>{\centering\arraybackslash}p{0.18\textwidth}
@{}}
\toprule
\textbf{Configuration}
&
\textbf{PPL}
&
\textbf{Absolute \(\Delta\)PPL}
&
\textbf{Relative \(\Delta\)PPL}
\\
\midrule

BF16 baseline
&
7.125503
&
0.000000
&
0.00\%
\\

SmoothQuant W8A8KV16
&
7.159544
&
+0.034041
&
+0.48\%
\\

NestQuant W4.25A16KV16
&
7.171211
&
+0.045708
&
+0.64\%
\\

AWQ W4A16KV16
&
7.418013
&
+0.292510
&
+4.11\%
\\

NestQuant W4.25A4KV4
&
7.421992
&
+0.296489
&
+4.16\%
\\

Project-specific W4A8
&
7.488315
&
+0.362812
&
+5.09\%
\\

\bottomrule
\end{tabularx}
\end{spacing}
\end{table}

Table~\ref{tab:app-memory-results} reports the runtime model-tensor
storage and peak GPU-memory measurements for the principal
cross-method comparison.

\begin{table}[htbp]
\centering
\caption{Complete runtime-storage and peak-GPU-memory results.}
\label{tab:app-memory-results}

\begin{spacing}{1}
\small
\renewcommand{\arraystretch}{1.18}
\setlength{\tabcolsep}{4pt}

\begin{tabularx}{\textwidth}{@{}
>{\raggedright\arraybackslash}X
>{\centering\arraybackslash}p{0.17\textwidth}
>{\centering\arraybackslash}p{0.17\textwidth}
>{\centering\arraybackslash}p{0.17\textwidth}
>{\centering\arraybackslash}p{0.17\textwidth}
@{}}
\toprule
\textbf{Configuration}
&
\textbf{Runtime storage (GiB)}
&
\textbf{Storage reduction}
&
\textbf{Peak GPU (GiB)}
&
\textbf{GPU reduction}
\\
\midrule

BF16
&
5.98
&
---
&
6.33
&
---
\\

SmoothQuant W8A8KV16
&
3.36
&
43.8\%
&
3.70
&
41.5\%
\\

AWQ W4A16KV16
&
2.13
&
64.4\%
&
2.47
&
61.0\%
\\

Project-specific W4A8
&
2.13
&
64.4\%
&
2.47
&
61.0\%
\\

\bottomrule
\end{tabularx}
\end{spacing}
\end{table}

The nearly identical storage and peak-memory results of AWQ and the
project-specific W4A8 pipeline indicate that their packed 4-bit weight
representations dominate the persistent memory reduction. Dynamic A8
processing changes the numerical execution path but does not materially
reduce the persistent model-tensor storage beyond the packed W4
representation.

Table~\ref{tab:app-nestquant-results} reports the complete NestQuant
quantization-scope ablation relative to W4.25A16KV16.

\begin{table}[htbp]
\centering
\caption{Complete NestQuant quantization-scope ablation results.}
\label{tab:app-nestquant-results}

\begin{spacing}{1}
\small
\renewcommand{\arraystretch}{1.18}

\begin{tabularx}{0.92\textwidth}{@{}
>{\raggedright\arraybackslash}X
>{\centering\arraybackslash}p{0.18\textwidth}
>{\centering\arraybackslash}p{0.20\textwidth}
>{\centering\arraybackslash}p{0.18\textwidth}
@{}}
\toprule
\textbf{Configuration}
&
\textbf{PPL}
&
\textbf{Absolute \(\Delta\)PPL}
&
\textbf{Relative \(\Delta\)PPL}
\\
\midrule

W4.25A16KV16
&
7.171211
&
0.000000
&
0.000\%
\\

W4.25A16KV4
&
7.224471
&
+0.053260
&
+0.743\%
\\

W4.25A4KV16
&
7.382130
&
+0.210919
&
+2.941\%
\\

W4.25A4KV4
&
7.421992
&
+0.250781
&
+3.497\%
\\

\bottomrule
\end{tabularx}
\end{spacing}
\end{table}

The ablation shows that activation quantization contributes the larger
quality penalty. The relative increase caused by A4 is approximately
four times the increase caused by KV4 under the retained configuration.
The combined A4+KV4 configuration produces the largest degradation but
also represents the broadest quantization scope considered in the
NestQuant quality experiments.

Finally, Table~\ref{tab:app-packed-nestquant} summarises the physical
deployment measurements obtained for packed NestQuant
W4.25A16KV16.

\begin{table}[htbp]
\centering
\caption{Packed NestQuant W4.25A16KV16 deployment results.}
\label{tab:app-packed-nestquant}

\begin{spacing}{1}
\small
\renewcommand{\arraystretch}{1.18}

\begin{tabularx}{0.86\textwidth}{@{}
>{\raggedright\arraybackslash}X
>{\centering\arraybackslash}p{0.21\textwidth}
>{\centering\arraybackslash}p{0.23\textwidth}
>{\centering\arraybackslash}p{0.19\textwidth}
@{}}
\toprule
\textbf{Metric}
&
\textbf{FP16 reference}
&
\textbf{Packed NestQuant}
&
\textbf{Reduction}
\\
\midrule

Checkpoint size
&
5.78 GiB
&
2.16 GiB
&
62.6\%
\\

Peak GPU memory
&
6.01 GiB
&
2.84 GiB
&
52.7\%
\\

\bottomrule
\end{tabularx}
\end{spacing}
\end{table}

These packed-deployment measurements demonstrate that the retained
W4.25 representation produces a substantially smaller physical model
artefact and a lower peak GPU-memory requirement. The corresponding
claim is limited to the packed W4.25A16KV16 deployment path; equivalent
packed A4 and KV4 deployment savings were not measured in the present
study.